\subsection{Extension of scalars in homomorphism modules}

Let A and B be two rings, E a right A-module, F a right B-module and G a (B, A)-bimodule. Recall that we have defined (Algebra, Chapter II, § 4, no. 2) a canonical homomorphism of Z-modules

\[
\nu \colon \mathbf {F} \otimes_ {\mathbf {B}} \operatorname{Hom} _ {\mathbf {A}} (\mathbf {E}, \mathbf {G}) \rightarrow \operatorname{Hom} _ {\mathbf {A}} (\mathbf {E}, \mathbf {F} \otimes_ {\mathbf {B}} \mathbf {G})\tag{8}
\]

such that, for all \(y \in \mathbf{F}\) and \(u \in \operatorname{Hom}_{\mathbf{A}}(\mathbf{E}, \mathbf{G})\) , \(\nu(y \otimes u)\) is the A-linear mapping \(x \mapsto y \otimes u(x)\) .

PROPOSITION 10. Let A, B be two rings, E a right A-module, F a right B-module and G a (B, A)-bimodule. Suppose that F is flat. Then, if E is finitely generated (resp. finitely presented), the canonical homomorphism (8) is injective (resp. bijective).

Consider A, B, F, G as fixed and for each right A-module E set

\[
\mathrm{T} (\mathrm{E}) = \mathrm{F} \otimes_ {\mathrm{B}} \operatorname{Hom} _ {\mathrm{A}} (\mathrm{E}, \mathrm{G}), \quad \mathrm{T} ^ {\prime} (\mathrm{E}) = \operatorname{Hom} _ {\mathrm{A}} (\mathrm{E}, \mathrm{F} \otimes_ {\mathrm{B}} \mathrm{G})
\]

and denote the homomorphism (8) by \(\nu_{\mathbf{E}}\) ; for every right A-module homomorphism \(\nu: \mathbf{E} \to \mathbf{E}'\) , set \(\mathrm{T}(v) = 1_{\mathbf{F}} \otimes \mathrm{Hom}(v, 1_{\mathbf{G}})\) and \(\mathrm{T}'(v) = \mathrm{Hom}(v, 1_{\mathbf{F}} \otimes 1_{\mathbf{G}})\) . Let \(\mathrm{L}_1 \xrightarrow{v} \mathrm{L}_0 \xrightarrow{w} \mathrm{E} \to 0\) be a presentation of E; we suppose the free module \(\mathrm{L}_0\) (resp. the free modules \(\mathrm{L}_0\) and \(\mathrm{L}_1\) ) to be finitely generated. The diagram

\[
\begin{array}{c} 0 \longrightarrow \mathrm{T(E)} \xrightarrow {\mathrm{T(w)}} \mathrm {T(L_ {0})} \xrightarrow {\mathrm{T(v)}} \mathrm {T(L_ {1})} \\ v _ {\mathbf {E}} \Biggl \downarrow \qquad \qquad \qquad \qquad \qquad \qquad \nu_ {\mathrm {L_ {0}}} \Biggl \downarrow \qquad \qquad \qquad \nu_ {\mathrm {L_ {1}}} \Biggl \downarrow \\ 0 \longrightarrow \mathrm {T^ {\prime} (E)} \xrightarrow {\mathrm {T^ {\prime} (w)}} \mathrm {T^ {\prime} (L_ {0})} \xrightarrow {\mathrm {T^ {\prime} (v)}} \mathrm {T^ {\prime} (L_ {1})} \end{array}\tag{9}
\]

is commutative and the second row is exact (Algebra, Chapter II, § 2, no. 1, Theorem 1); moreover, the sequence

\[
0 \rightarrow \operatorname{Hom} _ {\mathbf {A}} (\mathrm{E}, \mathrm{G}) \rightarrow \operatorname{Hom} _ {\mathbf {A}} (\mathrm{L} _ {0}, \mathrm{G}) \rightarrow \operatorname{Hom} _ {\mathbf {A}} (\mathrm{L} _ {1}, \mathrm{G})
\]

is exact (loc. cit.), and as F is flat, the first row of (9) is also an exact sequence (no. 3, Proposition 1). Then we know that \(\nu_{\mathrm{L}_0}\) (resp. \(\nu_{\mathrm{L}_0}\) and \(\nu_{\mathrm{L}_1}\) ) is bijective (resp. are bijective) (Algebra, Chapter II, § 4, no. 2, Proposition 2). If we assume only that \(\nu_{\mathrm{L}_0}\) is bijective, it follows from (9) that

\[
\nu_ {\mathbf {L} _ {0}} \circ \mathbf {T} (w) = \mathbf {T} ^ {\prime} (w) \circ \nu_ {\mathbf {E}}
\]

is injective and hence so is \(v_{E}\) . If we assume that both \(v_{L_{0}}\) and \(v_{L_{1}}\) are bijective, it follows from § 1, no. 4, Corollary 2, (ii) to Proposition 2 that \(v_{E}\) is surjective, and as we have just seen that \(v_{E}\) is injective, it is bijective.

